\section{Method}
\label{sec:method}

\subsection{Background for notation}
Given a video with RGB frames $\{\mathbf{I}^{\text{rgb}}_t\}_{t=1}^{T}$, our pipeline outputs a set of bird bounding boxes
$\mathcal{B}_t=\{b_t^n\}_{n=1}^{N_t}$ per frame and a set of tracks $\mathcal{T}_t=\{\tau_t^i\}_{i=1}^{M_t}$ over time.
We use pixel coordinates $\mathbf{x}=(x,y)^\top$.

\begin{figure*}[t]
  \centering
  \includegraphics[width=0.45\textwidth]{pic_framework.png}
  \caption{Birds' Track Framework.}
  \label{fig:birds_framework}
\end{figure*}

\paragraph{Preprocessing outputs.}
From each frame we compute:
(1) a smoothed grayscale intensity image $I_t \in [0,255]^{H\times W}$,
(2) a Laplacian-of-Gaussian (LoG) feature map $L_t \in [0,255]^{H\times W}$,
(3) a hard valid mask $V_t \in \{0,255\}^{H\times W}$ (e.g., subtitles as invalid), and
(4) a soft specular/glare mask $S_t \in \{0,255\}^{H\times W}$ (glare-like regions as $0$).

\paragraph{Camera motion model.}
We estimate an affine warp $\mathbf{T}_t\in\mathbb{R}^{2\times 3}$ that aligns frame $(t\!-\!1)$ to $t$.
Let $\mathcal{W}(\mathbf{x};\mathbf{T}_t)$ denote the warped coordinate. The aligned previous feature is
\begin{equation}
\widehat{X}_{t-1\rightarrow t}(\mathbf{x}) = X_{t-1}\big(\mathcal{W}(\mathbf{x};\mathbf{T}_t)\big),
\end{equation}
where $X\in\{I,L\}$ depending on the stage.

\paragraph{Track state.}
Each track $\tau^i$ maintains a 2D bounding-box state
\begin{equation}
\mathbf{s}^i_t = [c^i_{x,t},\,c^i_{y,t},\,w^i_t,\,h^i_t]^\top,
\end{equation}
where $(c_x,c_y)$ denotes the box center and $(w,h)$ its width and height.
The corresponding bounding box is decoded as
\[
b(\mathbf{s}^i_t)=\big(\lfloor c_x-\tfrac{w}{2}\rfloor,\ \lfloor c_y-\tfrac{h}{2}\rfloor,\ w,\ h\big).
\]

\subsection{Preprocessing}
\label{sec:preprocessing}

Bird flocks in natural scenes often appear with strong illumination changes, glare, and embedded subtitles. Subtitles create hard edges that severely interfere with both camera motion estimation and bird motion detection. To address this, we compute two complementary feature streams along with specialized masks.

\paragraph{Subtitle Masking.}
We detect subtitle regions using a fixed-parameter specular-like detector applied only to the bottom 26\% of the image (empirical subtitle region). The detector identifies high-value pixels with low saturation in HSV space and strong local contrast (via top-hat filtering: $V - G_\sigma * V$). We then apply a two-stage morphological dilation: (1) strong horizontal dilation (25$\times$3) to group characters into line-level blocks for connected component filtering, and (2) light dilation (3$\times$1) to fill broken strokes without bridging inter-character gaps (avoid hidding the birds which is in the subtitle gap). Components are filtered by area, width, height, and vertical position, yielding a hard binary mask where subtitle regions are marked as invalid (0). These regions are excluded from all subsequent motion computations.

\paragraph{Mask Hole Filling.}
Invalid subtitle pixels are replaced by Gaussian-blurred ($\sigma=3$) intensity values to avoid creating artificial edges that disrupt phase correlation and differencing operations.

\paragraph{Specular Mask.} 
Glare from water surfaces or direct sunlight creates false motion signals. However, some bird species exhibit specular-like appearances. Our soft glare mask distinguishes problematic glare from bird specularity. We detect glare pixels using adaptive thresholds based on per-video brightness statistics (99th percentile of $V$ channel). Detected seeds are grown into coherent regions using brightness-constrained dilation. Connected components are filtered by area (5px to 5\% of image) and local texture density (edge density $> 0.01$ to avoid smooth surfaces like sky). A fallback mechanism disables glare suppression when total glare area exceeds 6\% or when a single component dominates, preventing over-suppression of legitimate bird motion.

\paragraph{Background Smoothing and LoG Feature.}
We apply bilateral filtering ($d=7$, $\sigma_{\text{color}}=35$, $\sigma_{\text{space}}=15$) to suppress high-frequency texture while preserving edges (many motion pixels are on the egdes). The smoothed grayscale serves as the intensity feature. For robustness to global illumination changes, we compute a Laplacian-of-Gaussian (LoG) feature by applying a Laplacian kernel ($k=3$) to Gaussian-blurred ($\sigma=1.2$) intensity, followed by 1–99 percentile normalization to 0–255 range.

\subsection{Camera Motion Compensation}
\label{sec:camera_motion_comp}

Many videos exhibit subtle camera motion (e.g., from wind or tripod adjustment) that is imperceptible to human eyes but detectable through stationary background features like distant mountains. Over-compensation for such motion would harm bird detection, while under-compensation allows camera motion to dominate differential signals. We therefore adopt a conservative translation-only model (not Euclidean or full affine) with loose acceptance criteria.

\paragraph{Three-Stage Estimation Pipeline.}
We employ a cascaded approach that progressively refines the estimate:
\begin{enumerate}
    \item \textbf{Global Phase Correlation:} Compute translation between full frames using phase correlation with optional Hanning window. Accept if response $\geq 0.15$ and alignment error improves by $>5\%$.
    \item \textbf{ROI Consensus Refinement:} If global estimate is unreliable, partition image into ROIs (corners and/or strips). For each ROI, compute phase correlation shift weighted by response, valid pixel fraction, and Shi–Tomasi corner count. Compute weighted median shift across ROIs, reject outliers via MAD (median absolute deviation), and re-compute weighted median from inliers.
    \item \textbf{ECC Fallback:} If ROI consensus fails, use Enhanced Correlation Coefficient (ECC) maximization with translation motion model, initialized by previous best estimate.
\end{enumerate}

\paragraph{Motion Decision and Warping.}
The final translation $(dx, dy)$ is clamped to $\pm 50$ pixels. Camera is considered "moving" if $\sqrt{dx^2 + dy^2} \geq 0.5$ pixels. The previous frame is warped using the translation matrix $\mathbf{T} = \begin{bmatrix} 1 & 0 & dx \\ 0 & 1 & dy \end{bmatrix}$ with BORDER\_REPLICATE interpolation.

\paragraph{Specular-Aware Error Metric.}
When evaluating alignment quality, we down-weight errors inside specular regions by a factor of 0.6 (soft weighting) to prevent glare from dominating the error metric.

\subsection{Motion-Based Candidate Generation}
\label{sec:candidate_gen}

Given aligned intensity features $I_t$ and $\hat{I}_{t-1 \rightarrow t}$, we generate bird motion candidates through a multi-branch system designed for high recall in diverse scenarios.

\paragraph{Core Difference Map.}
We compute a normalized intensity difference map:
\[
D^{\text{norm}}_t(\mathbf{x}) = \frac{|I_t(\mathbf{x}) - \hat{I}_{t-1 \rightarrow t}(\mathbf{x})|}{I_t(\mathbf{x}) + \hat{I}_{t-1 \rightarrow t}(\mathbf{x}) + 3.0} \times 255
\]
which is robust to global brightness scaling. A Gaussian blur (kernel size 5) suppresses glare noise.

\paragraph{Adaptive Thresholding.}
We compute a robust threshold via MAD: $\text{threshold} = \text{median}(D) + 8 \times 1.4826 \times \text{MAD}$, clamped between 8 and 80. This adapts to varying motion magnitudes across videos.

\paragraph{Multi-Branch Enhancement.}
\begin{itemize}
    \item \textbf{LoG Difference Branch:} Automatically enabled when edge density $> 0.06$ or glare/specular ratio exceeds thresholds. Computes absolute LoG difference with separate MAD threshold ($k=7$). Fused with intensity difference via OR (normal scenes) or AND (complex/flickery scenes).
    \item \textbf{KNN Background Subtraction:} Used primarily for \textit{static camera} scenarios where birds move slowly against stable backgrounds. Automatically enabled when camera motion EMA $< 0.25$ and foreground ratio $> 0.10$. Learning rate is reduced to 0 when foreground is high to prevent birds from being absorbed into background.
\end{itemize}

\paragraph{Background-Only Component Suppression.}
A key innovation addresses dynamic backgrounds (e.g., waves, foliage motion). Components detected mainly by KNN but with low support from core difference map (overlap $< 0.06$) are considered "background-only." We apply two gates: (1) mean difference must be at least 2.0 above threshold, and (2) after 5-frame warmup, temporal persistence score must exceed 1.6 (equivalent to $\geq$ 2 consecutive frames). This effectively suppresses repetitive background motion while preserving true birds.

\paragraph{Morphology and Filtering.}
We apply resolution-adaptive morphology: opening (kernel size 3 at 720p) removes noise, closing (size 9) fills gaps, and optional bridge operation connects broken bird parts (the not motion pixels in the birds cause broken) . Connected components are filtered by area (80–infinity pixels), aspect ratio ($\leq 10$), and minimum dimension (5 pixels). Nested boxes are suppressed using IoA (Intersection over Area) threshold of 0.92.

\paragraph{Temporal Scoring.}
Components are scored by persistence (exponential decay with factor 0.85) and edge change (mean LoG difference within box). Top 12 boxes by combined score are kept as final candidates.


\subsection{Candidate Refinement}
The motion mask can still contain background fragments (waves, foliage). We refine each raw box using ROI-level instance separation and birdness constraints.

\paragraph{ROI foreground cleanup.}
For each raw box $b=(x,y,w,h)$, we crop the ROI from $M_t$ and apply open/close. If the number of foreground pixels is too small, we discard it.

\paragraph{Watershed-based splitting.}
Let $\text{fg}\in\{0,1\}^{h\times w}$ be the ROI foreground. We compute a distance transform:
\begin{equation}
\text{dist}=\texttt{DT}(\text{fg}),
\end{equation}
and build sure-foreground markers by thresholding:
\begin{equation}
\text{sure\_fg}=\mathbb{I}\left[\text{dist} > \rho\cdot \max(\text{dist})\right],\quad \rho\in(0,1).
\end{equation}
We then run watershed on the RGB ROI to separate instances and convert each instance into a bounding box.

\paragraph{Geometric and photometric constraints.}
We keep an instance box only if it satisfies:
(i) area bounds relative to the full image,
(ii) minimum side length,
(iii) aspect ratio bounds,
(iv) foreground occupancy inside the box (foreground ratio),
(v) instance fill ratio,
and (vi) a glare constraint (discard boxes mostly inside spec regions).
Per raw box, we keep the largest remaining instance to avoid duplicated boxes for one bird.

\subsection{Tracking}
We associate refined boxes across time and output stable tracklets. Our method uses: (i) IoU-gated detection correction,
(ii) IoA-based birth suppression and duplicate removal, (iii) simple track management with age and hit counters, and (iv) Iterative Lucas--Kanade (IK-LK) refinement (as specified in the technical report).

\paragraph{Detection correction with IoU gating.}
After LK propagation, we associate detections to existing tracks using
Intersection-over-Union (IoU), defined as
\begin{equation}
\text{IoU}(b^i, b^j) = \frac{|b^i \cap b^j|}{|b^i \cup b^j|}.
\end{equation}
For each track, we select the unmatched detection with the highest IoU.
If the IoU exceeds a threshold $\tau_{\text{iou}}$, the track is corrected by
replacing its bounding box with the detection box, and the track hit counter is
incremented.
This correction step prevents long-term drift caused by optical flow errors.

\paragraph{New track initialization with IoA suppression.}
Unmatched detections may spawn new tracks.
To avoid duplicate tracks caused by fragmented or nested detections,
we apply an Intersection-over-Area (IoA) suppression rule.
IoA is defined as
\begin{equation}
\text{IoA}(b_{\text{inner}}, b_{\text{outer}})
= \frac{|b_{\text{inner}} \cap b_{\text{outer}}|}{|b_{\text{inner}}|}.
\end{equation}
A detection box $b^{\text{new}}$ is discarded if it is largely contained within
any existing track box $b^i$:
\begin{equation}
\text{IoA}(b^{\text{new}}, b^i) > \tau_{\text{ioa}}.
\end{equation}
Only detections that pass this test are used to initialize new tracks.

\paragraph{Track management.}
Each track maintains a hit counter and a miss counter.
Unmatched tracks increase a miss counter and are removed if
$\text{miss}>\text{max\_age}$.
A track is reported as \emph{confirmed} only after reaching
$\text{min\_hits}$ matches to suppress flickering false positives.

Unmatched detections spawn new tracks only if they are sufficiently
separated from existing tracks (birth gate):
\begin{equation}
\max_i \text{IoU}(b^i,b^{\text{new}}) \le \tau_{\text{birth-iou}}
\quad\wedge\quad
\min_i \frac{\|\mathbf{c}(b^i)-\mathbf{c}(b^{\text{new}})\|_2}{\text{diag}(b^i)}
\ge \tau_{\text{birth-dist}}.
\end{equation}
This gate prevents creating duplicate tracks for nearby detections.

\paragraph{Iterative Lucas--Kanade Algorithm (general formulation, iterative version).}
To improve short-term localization and robustness to missed detections, we refine per-track displacement between consecutive frames using IK-LK.
Let $I(\mathbf{x})$ be the template patch from the previous frame (e.g., around the predicted track location),
and let $J(\mathbf{x})$ be the current-frame patch. We seek a translation $\mathbf{d}$ such that
$I(\mathbf{x}) \approx J(\mathbf{x}+\mathbf{d})$.
We initialize $\bar{\nu}_0$ by the motion-model prediction (or zero), and iterate:

\begin{align}
\bar{\eta}_k &= G^{-1}\,b_k,\\
b_k &= \sum_{\mathbf{x}}\begin{bmatrix}
\delta I_k(\mathbf{x})\,I_x(\mathbf{x})\\
\delta I_k(\mathbf{x})\,I_y(\mathbf{x})
\end{bmatrix},\qquad
\delta I_k(\mathbf{x}) = I(\mathbf{x}) - J_k(\mathbf{x}),\\
J_k(\mathbf{x}) &= J(\mathbf{x}+\bar{\nu}_{k-1}),
\end{align}
where $I_x,I_y$ are spatial gradients of the template (computed once), and
\begin{equation}
G=\sum_{\mathbf{x}}
\begin{bmatrix}
I_x(\mathbf{x})^2 & I_x(\mathbf{x})I_y(\mathbf{x})\\
I_x(\mathbf{x})I_y(\mathbf{x}) & I_y(\mathbf{x})^2
\end{bmatrix}
\end{equation}
is constant across iterations. After computing $\bar{\eta}_k$, we update the displacement hypothesis:
\begin{equation}
\bar{\nu}_k=\bar{\nu}_{k-1}+\bar{\eta}_k.
\end{equation}
We iterate until $\|\bar{\eta}_k\|_2$ is below a threshold or a fixed number of iterations is reached. The final displacement is
\begin{equation}
\bar{\nu}=\mathbf{d}=\sum_{k=1}^{K}\bar{\eta}_k=\bar{\nu}_K.
\end{equation}

\paragraph{Using IK-LK in tracking.}
For each active track, we run IK-LK on an intensity patch (optionally after camera-motion compensation) and refine the predicted center:
$\mathbf{c}^{i-}_t \leftarrow \mathbf{c}^{i-}_t + \mathbf{d}$.
This refined prediction is then used for association (or, if no detection matches, for short-term propagation within $\text{max\_age}$).